% Part: incompleteness
% Chapter: theories-computability
% Section: complete-decidable

\documentclass[../../../include/open-logic-section]{subfiles}

\begin{document}

\olfileid[hi]{inc}{tcp}{cdc}

\olsection{\printtoken{S}{axiomatizable} पूर्ण सिद्धांत निर्णेय हैं}

कोई सिद्धांत {\em पूर्ण} कहलाता है यदि प्रत्येक वाक्य $!A$ के लिए या तो
$!A$ अथवा $\lnot !A$ सिद्ध हो।

\begin{lem}
मान लें कोई सिद्धांत $\Th{T}$ पूर्ण और !!{axiomatizable} है। तब $\Th{T}$
निर्णेय है।
\end{lem}

\begin{proof}
मान लें $\Th{T}$ पूर्ण है और $A$ अभिगृहीतों का संगणनीय समुच्चय है। यदि
$\Th{T}$ असंगत है, तो स्पष्टतः वह संगणनीय है। (कलन-विधि: ``बस हाँ कहें।'')
अतः हम मान सकते हैं कि $\Th{T}$ संगत भी है।

यह तय करने के लिए कि कोई वाक्य $!A$,~$\Th{T}$ में है या नहीं,~$\Th{T}$
से~$!A$ की किसी !!a{derivation} और~$\lnot !A$ की किसी !!a{derivation} की
एक साथ खोज करें। $\Th{T}$ पूर्ण है, इसलिए दोनों में से कोई एक अवश्य
मिलेगी; और $\Th{T}$ संगत है, इसलिए यदि आपको~$\lnot !A$ की
!!a{derivation} मिले, तो~$!A$ की कोई !!{derivation} नहीं है।

भिन्न शब्दों में, हम पहले ही जानते हैं कि $\Th{T}$,~!!{c.e.} है; इसलिए
पहले सिद्ध किए गए एक प्रमेय के अनुसार इतना दिखाना पर्याप्त है कि
$\Th{T}$ का पूरक भी~!!{c.e.} है। किंतु कोई !!a{formula}~$!A$,
$\Th{\bar T}$ में तभी और केवल तभी है जब $\lnot !A$,~$\Th{T}$ में हो;
अतः $\Th{\bar T} \leq_m \Th{T}$।
\end{proof}

\end{document}
